\SourceLeaf{215}%
\[
\begin{aligned}
\frac{\mathrm{d}m}{m}&=\mathrm{cotang}\,U\mathord{.}\mathrm{d}u-\mathrm{cotang}\,w\mathord{.}\mathrm{d}w+\frac{\varepsilon\varepsilon\cos w\mathord{.}\sin w\mathord{.}\mathrm{d}w}{1-\varepsilon\varepsilon\cos w^2}\\
\frac{\mathrm{d}U}{\sin U}&=\frac{\mathrm{d}w}{\sin w}-\frac{\varepsilon\varepsilon\sin w\mathord{.}\mathrm{d}w}{1-\varepsilon\varepsilon\cos w^2}=\frac{(1-\varepsilon\varepsilon)\mathrm{d}w}{(1-\varepsilon\varepsilon\cos w^2)\sin w}
\end{aligned}
\]
\noindent und hieraus
\[
\frac{\mathrm{d}m}{m}=\frac{(1-\varepsilon\varepsilon)\mathrm{d}w}{\sin w(1-\varepsilon\varepsilon\cos w^2)}\mathord{.}(\cos U-\cos w)
\]
\noindent Hieraus erhellt, dass \(m\) da seinen grössten oder kleinsten Werth erhält, wo \(U=w\) wird; bezeichnet man den Werth von \(w\) an dieser Stelle durch \(W\), so wird
\[
k=(\frac{1-\varepsilon\cos W}{1+\varepsilon\cos W})^{\frac{1}{2}\varepsilon}\quad\text{oder}\quad \cos W=\frac{1-k^{\frac{2}{\varepsilon}}}{\varepsilon(1+k^{\frac{2}{\varepsilon}})}
\]
\noindent woraus man \(W\) bestimmen kann, wenn \(k\) nach der obigen Formel berechnet ist. Für die Ausübung wird inzwischen auf die ganz genaue Gleichheit der Werthe von \(m\) an den äussersten Breiten wenig ankommen, und man kann sich begnügen, für \(90^0-W\) ungefähr die mittlere Breite zu wählen, und daraus \(k\) abzuleiten. Den allgemeinen Zusammenhang zwischen \(U\) und \(w\) gibt dann die Formel
\[
\mathrm{tang}\tfrac{1}{2}U=\mathrm{tang}\tfrac{1}{2}w\left\{\frac{(1-\varepsilon\cos W)(1+\varepsilon\cos w)}{(1+\varepsilon\cos W)(1-\varepsilon\cos w)}\right\}^{\frac{1}{2}\varepsilon}
\]
\noindent Zur wirklichen numerischen Berechnung ist es jedoch vortheilhafter, Reihen anzuwenden, denen man verschiedene Formen geben kann, bei deren Entwicklung wir uns aber hier nicht aufhalten.

Da man übrigens leicht sieht, dass für \(w<W\), \(U>w\), also \(\cos U-\cos w\) und mithin auch \(\frac{\mathrm{d}m}{\mathrm{d}w}\) negativ; und für \(w>W\), \(U<w\), mithin \(\frac{\mathrm{d}m}{\mathrm{d}w}\) positiv wird, so ist klar, dass für \(w=U=W\) der Werth von \(m\) allemal ein Minimum wird, und zwar
\[
=\frac{A}{a}\BareRoot{(1-\varepsilon\varepsilon\cos W^2)}
\]
\noindent Wählt man also den Halbmesser der Kugel \(A=\frac{a}{\BareRoot{(1-\varepsilon\varepsilon\cos W^2)}}\), so ist die Darstellung unendlich kleiner Theile des Ellipsoids bei der Breite \(90^0-W\) dem Urbilde nicht bloss ähnlich, sondern gleich, bei andern Breiten aber grösser.

Man kann den Logarithmen von \(m\) mit Vortheil in eine nach den Potenzen von \(\cos U-\cos W\) fortlaufende Reihe entwickeln, deren erste für die Ausübung zureichende Glieder diese sind
